\documentclass[a4paper,10pt]{article}

% Packages
\usepackage{txfonts}
\usepackage{titlesec}
\usepackage[utf8]{inputenc}
\usepackage[T1]{fontenc}
\usepackage[french]{babel}
\usepackage{geometry}
\usepackage{graphicx}
\usepackage[intoc]{nomencl}
\usepackage{multicol}
\usepackage{fancyhdr}
\usepackage{hyperref}
\usepackage{float}
\usepackage{setspace}
\usepackage{caption}
\usepackage{booktabs}
\usepackage{amsmath}
\usepackage{siunitx}
\usepackage{textcomp} % pour \textmu

% Configuration siunitx
\sisetup{
  locale = FR,
  table-number-alignment = center
}

% Mise en page
\geometry{a4paper, top=18mm, bottom=18mm, textwidth=185mm}
\setlength{\columnsep}{6mm}
\setlength{\parindent}{0pt}

\titleformat{\section}
{\normalfont\bfseries}
{\thesection.}{0.5em}{}

\titleformat{\subsection}
{\normalfont\itshape}
{\thesubsection.}{0.5em}{}

\pagestyle{fancy}
\renewcommand{\footrulewidth}{0pt}
\fancyhf{}
\rhead{\includegraphics[scale=0.4]{Images/hepia_simple.png}}
\rfoot{\textit{\today}}
\cfoot{\thepage}

\captionsetup{
  font=footnotesize,
  labelsep=colon
}

% Nomenclature
\makenomenclature
\nomenclature{AFM}{Microscope à Force Atomique}
\nomenclature{$R_a$}{Rugosité moyenne arithmétique}
\nomenclature{$R_q$}{Rugosité RMS (écart-type)}
\nomenclature{$R_{pv}$}{Peak-to-valley (écart max--min)}
\nomenclature{PSPD}{Photodiode sensible à la position (4 quadrants)}
\nomenclature{LFM}{Microscopie en force latérale (friction)}

\begin{document}

\begin{center}
\vspace*{0mm}

\large
Rapport de laboratoire\\
\vspace{3mm}

SIMU ELECTROMAG\\
\vspace{8mm}

Guilhem Rozier Vilardell\footnote{\textit{\href{mailto:guilhem.roziervi@hes-so.ch}{MT2, guilhem.roziervi@hes-so.ch}}},

Noah Michelot\footnote{\textit{\href{mailto:noah.michelot@hes-so.ch}{MT2, noah.michelot@hes-so.ch}}},

Mattéo Pozzo Di Borgo\footnote{\textit{\href{mailto:matteo.pozzodi@hes-so.ch}{MT2, matteo.pozzodi@hes-so.ch}}}

\vspace{3mm}

\footnotesize{
\textit{Haute école de paysage, d'ingénierie et d'architecture, Genève, Suisse}
}\\

\vspace{9mm}
\end{center}

\rule{\linewidth}{0.4pt}

\onehalfspacing

% ============================================================
% TABLEAU 1
% ============================================================

\begin{table}[H]
\centering

\caption{
Mesure de l'intensité des tensions $U_{z,\mathrm{eff}}$ et
$U_{x,\mathrm{eff}}$ en fonction de la position axiale $z$.
$f = 900\,\mathrm{Hz}$,
$I_\mathrm{eff} = 0.9\,\mathrm{A}$.
}

\label{tab:comparaison_B}

\small

\begin{tabular}{
S[table-format=3.0]
S[table-format=2.2]
S[table-format=1.2]
}

\toprule

{$z \pm 1$ / mm} &
{$U_{z,\mathrm{eff}} \pm 0.01$ / mV} &
{$U_{x,\mathrm{eff}} \pm 0.01$ / mV} \\

\midrule

0   & 42.66 & 0.79 \\
5   & 42.29 & 0.81 \\
10  & 41.57 & 0.93 \\
15  & 40.24 & 0.77 \\
20  & 38.61 & 0.72 \\
25  & 36.74 & 0.66 \\
30  & 34.76 & 0.65 \\
35  & 32.55 & 0.53 \\
40  & 30.34 & 0.47 \\
45  & 28.16 & 0.42 \\
50  & 25.99 & 0.41 \\
55  & 23.73 & 0.41 \\
60  & 21.89 & 0.41 \\
65  & 19.98 & 0.41 \\
75  & 16.74 & 0.41 \\
85  & 13.90 & 0.41 \\
95  & 11.64 & 0.41 \\
105 & 9.69  & 0.41 \\
115 & 8.22  & 0.41 \\
125 & 6.91  & 0.41 \\
145 & 5.01  & 0.41 \\
165 & 3.72  & 0.41 \\
185 & 2.83  & 0.41 \\
205 & 2.20  & 0.41 \\
225 & 1.73  & 0.41 \\
245 & 1.40  & 0.41 \\
265 & 1.14  & 0.41 \\
285 & 0.95  & 0.41 \\
305 & 0.80  & 0.41 \\
325 & 0.68  & 0.41 \\

\bottomrule

\end{tabular}

\end{table}

\vspace{5mm}

% ============================================================
% TABLEAU 2
% ============================================================

\begin{table}[H]
\centering

\caption{
Comparaison des valeurs de $B_\mathrm{eff}$ mesurées,
théoriques (MATLAB) et simulées pour quatre points de mesure.
$f = 900\,\mathrm{Hz}$,
$I_\mathrm{eff} = 0.9\,\mathrm{A}$.
}

\label{tab:Beff_compare_full}

\small

\begin{tabular}{
S[table-format=1.1]
S[table-format=1.1]
S[table-format=2.2]
S[table-format=1.2]
S[table-format=3.0]
S[table-format=2.1]
S[table-format=3.0]
S[table-format=3.0]
}

\toprule

{$z \pm 1.4$ / mm} &
{$x \pm 1.4$ / mm} &
{$U_{z,\mathrm{eff}} \pm 0.01$ / mV} &
{$U_{x,\mathrm{eff}} \pm 0.01$ / mV} &
{$B_\mathrm{eff,exp}$ / \textmu T} &
{$\sigma_{B_\mathrm{exp}}$ / \textmu T} &
{$B_\mathrm{eff,th}$ / \textmu T} &
{$B_\mathrm{simu}$ / \textmu T} \\

\midrule

67 & 198 & 12.80 & 9.65 & 247 & 18 & 264 & 264 \\
67 & 102 & 2.96  & 8.79 & 143 & 11 & 157 & 156 \\
107 & 102 & 2.88  & 4.18 & 78  & 6  & 84  & 84  \\
107 & 198 & 6.46  & 4.10 & 118 & 9  & 123 & 123 \\

\bottomrule

\end{tabular}

\smallskip

\raggedright
\footnotesize

$\varepsilon =
\lvert B_\mathrm{val} - B_\mathrm{exp} \rvert
/ B_\mathrm{exp} \times 100\,\%$.

Incertitude sur $B_\mathrm{exp}$ estimée à $\pm 5\,\%$
(propagation depuis
$\delta U = \pm 0.5\,\mathrm{mV}$
et
$\delta I = \pm 0.01\,\mathrm{A}$).

\end{table}  % <-- FIX: \end{table} manquant dans l'original

\vspace{5mm}

% ============================================================
% TABLEAU 2b (version courte)
% ============================================================
\begin{table}[H]
\centering

\caption{
Comparaison des valeurs de $B_\mathrm{eff}$ mesurées,
théoriques (MATLAB) et simulées pour quatre points de mesure.
$f = 900\,\mathrm{Hz}$,
$I_\mathrm{eff} = 0.9\,\mathrm{A}$.
}

\label{tab:Beff_compare_full}

\small

\begin{tabular}{
S[table-format=1.1]
S[table-format=1.1]
S[table-format=3.0]
S[table-format=2.1]
S[table-format=3.0]
S[table-format=3.0]
S[table-format=2.1]
}

\toprule

{$z \pm 1.4$ / mm} &
{$x \pm 1.4$ / mm} &
{$B_\mathrm{eff,exp}$ / \textmu T} &
{$\sigma_{B_\mathrm{exp}}$ / \textmu T} &
{$B_\mathrm{eff,th}$ / \textmu T} &
{$B_\mathrm{simu}$ / \textmu T} &
{$\varepsilon_{r,\mathrm{exp/simu}}$ / \%} \\

\midrule

67  & 198 & 247 & 18 & 264 & 264 & 6.5 \\
67  & 102 & 143 & 11 & 157 & 156 & 9.1 \\
107 & 102 &  78 &  6 &  84 &  84 & 7.7 \\
107 & 198 & 118 &  9 & 123 & 123 & 4.2 \\

\bottomrule

\end{tabular}

\smallskip
\raggedright
\footnotesize
$\varepsilon_{r,\mathrm{exp/simu}} =
\lvert B_\mathrm{simu} - B_\mathrm{eff,exp} \rvert / B_\mathrm{eff,exp} \times 100\,\%$.

\end{table}

\vspace{5mm}


\begin{table}[H]
\centering

\caption{
Comparaison des valeurs de $B_\mathrm{eff}$ mesurées,
théoriques (MATLAB) et simulées pour quatre points de mesure.
$f = 900\,\mathrm{Hz}$,
$I_\mathrm{eff} = 0.9\,\mathrm{A}$.
}

\label{tab:Beff_compare_full}

\small

\begin{tabular}{
S[table-format=3.2]
S[table-format=2.2]
S[table-format=3.2]
S[table-format=3.2]
S[table-format=2.2]
S[table-format=2.2]
}

\toprule

{$B_\mathrm{eff,exp}$ / \textmu T} &
{$\sigma_{B_\mathrm{exp}}$ / \textmu T} &
{$B_\mathrm{eff,th}$ / \textmu T} &
{$B_\mathrm{simu}$ / \textmu T} &
{$\varepsilon_{r,\mathrm{th}}$ / \%} &
{$\varepsilon_{r,\mathrm{simu}}$ / \%} \\

\midrule

247.49 & 18.33 & 263.58 & 264.34 & 6.50 & 6.81 \\
143.24 & 10.61 & 156.85 & 155.70 & 9.50 & 8.70 \\
 78.43 &  5.81 &  84.22 &  83.66 & 7.38 & 6.67 \\
118.18 &  8.75 & 123.29 & 123.28 & 4.32 & 4.32 \\

\bottomrule

\end{tabular}

\smallskip
\raggedright
\footnotesize
$\varepsilon_{r} =
\lvert B_\mathrm{val} - B_\mathrm{eff,exp} \rvert / B_\mathrm{eff,exp} \times 100\,\%$.

\end{table}



% ============================================================
% TABLEAU 3
% ============================================================

\begin{table}[H]
\centering

\caption{
Comparaison des valeurs de $B_\mathrm{eff}$ de la simulation COMSOL
en fonction de différents paramètres.
Les écarts relatifs sont calculés par rapport à
$B_\mathrm{simu,Comsol}$ :
\[
\varepsilon_r =
\frac{\left|B_\mathrm{val} - B_\mathrm{simu,Comsol}\right|}
{B_\mathrm{simu,Comsol}}
\times 100
\]
}

\label{tab:comsol}

\small

\begin{tabular}{
S[table-format=3.0]
S[table-format=3.0]
S[table-format=2.2]
S[table-format=3.0]
S[table-format=2.2]
S[table-format=3.0]
S[table-format=2.2]
}

\toprule

{$B_\mathrm{simu,Comsol}$ / \textmu T} &
{$B_\mathrm{simu,air,petit}$ / \textmu T} &
{$\varepsilon_{r,\mathrm{petit}}$ / \%} &
{$B_\mathrm{simu,air,grand}$ / \textmu T} &
{$\varepsilon_{r,\mathrm{grand}}$ / \%} &
{$B_\mathrm{simu,maillage}$ / \textmu T} &
{$\varepsilon_{r,\mathrm{maillage}}$ / \%} \\

\midrule

264 & 255 & 3.41 & 264 & 0.00 & 252 & 4.55 \\
156 & 154 & 1.28 & 154 & 1.28 & 171 & 9.62 \\
84  & 78  & 7.14 & 84  & 0.00 & 99  & 17.86 \\
123 & 113 & 8.13 & 122 & 0.81 & 135 & 9.76 \\

\bottomrule

\end{tabular}

\end{table}

% ============================================================
% SECTION : TENSION INDUITE
% ============================================================

\section{Partie tension induite}

\begin{table}[H]
  \centering
  \caption{Tension induite en fonction de la fréquence}
  \begin{tabular}{
    S[table-format=4.0]
    S[table-format=4.0]
    S[table-format=1.4]
    S[table-format=1.4]
    S[table-format=1.4]
    S[table-format=1.4]
    S[table-format=1.4]
  }

    \toprule

    {$f$ / Hz} &
    {$\omega$ / rad\,s$^{-1}$} &
    {$U_{i,30} \pm 0.0001$ / V} &
    {$U_{i,40} \pm 0.0001$ / V} &
    {$U_{i,55} \pm 0.0001$ / V} &
    {$U_{i,70} \pm 0.0001$ / V} &
    {$U_{i,100} \pm 0.0001$ / V} \\

    \midrule

    1000 & 6283 & 6.5461 & 5.1894 & 3.7373 & 3.1101 & 1.8441 \\
     900 & 5655 & 5.8753 & 4.7054 & 3.3818 & 2.5273 & 1.4926 \\
     800 & 5027 & 5.2405 & 4.1736 & 2.9970 & 2.2451 & 1.3398 \\
     700 & 4398 & 4.5708 & 3.6381 & 2.6304 & 1.9522 & 1.1704 \\
     600 & 3770 & 3.9461 & 3.1262 & 2.2515 & 1.6755 & 1.0080 \\
     500 & 3142 & 3.2802 & 2.6094 & 1.8777 & 1.4026 & 0.8348 \\
     400 & 2513 & 2.6222 & 2.0843 & 1.5026 & 1.1250 & 0.6674 \\
     300 & 1885 & 1.9728 & 1.5677 & 1.1265 & 0.8405 & 0.4985 \\

    \bottomrule

  \end{tabular}
\end{table}

% Tableau : mutuelle inductance brute
\begin{table}[H]
  \centering
  \caption{Mutuelle inductance en fonction de la distance}
  \begin{tabular}{
    S[table-format=3.0]
    S[table-format=1.2]
    S[table-format=2.1]
    S[table-format=3.0]
    S[table-format=1.2]
  }

    \toprule

    {$d \pm 1$ / mm} &
    {$a$ / mV\,s\,rad$^{-1}$} &
    {$\sigma_a$ / \textmu V\,rad$^{-1}$} &
    {$M_\mathrm{exp}$ / nH} &
    {$\sigma_{M_\mathrm{exp}}$ / nH} \\

    \midrule

    30  & 1.0  &  5.7 & 120 & 1.5  \\
    40  & 0.83 &  6.5 &  96 & 1.3  \\
    55  & 0.60 &  3.4 &  69 & 0.86 \\
    70  & 0.49 & 54.0 &  56 & 6.3  \\
    100 & 0.29 & 31.0 &  33 & 3.6  \\

    \bottomrule

  \end{tabular}
\end{table}

% Tableau : comparaison mutuelle inductance
\begin{table}[H]
  \centering
  \caption{Comparaison des mutuelles inductances mesurée, théorique et simulée}
  \begin{tabular}{
    S[table-format=3.0]
    S[table-format=3.0]
    S[table-format=1.2]
    S[table-format=3.0]
    S[table-format=2.1]
    S[table-format=3.0]
    S[table-format=2.1]
  }

    \toprule

    {$d \pm 1$ / mm} &
    {$M_\mathrm{exp}$ / nH} &
    {$\sigma_{M_\mathrm{exp}}$ / nH} &
    {$M_\mathrm{th}$ / nH} &
    {$\delta_\mathrm{th}$ / \%} \\

    \midrule

    30  & 120 & 1.5  & 118 &  1.7\\
    40  &  96 & 1.3  &  93 &  3.2 \\
    55  &  69 & 0.86 &  67 &  3.0  \\
    70  &  56 & 6.3  &  50 & 12.0  \\
    100 &  33 & 3.6  &  30 & 10.0 \\

    \bottomrule

  \end{tabular}
\end{table}


\begin{table}[H]
\centering

\caption{
Comparaison des mutuelles inductances mesurée, théorique et simulée.
$I_\mathrm{eff} = 0.9\,\mathrm{A}$.
}

\label{tab:mutual_compare}

\small

\begin{tabular}{
S[table-format=3.2]
S[table-format=1.2]
S[table-format=3.2]
S[table-format=3.2]
S[table-format=2.2]
S[table-format=2.2]
}

\toprule

{$M_\mathrm{exp}$ / nH} &
{$\sigma_{M_\mathrm{exp}}$ / nH} &
{$M_\mathrm{th}$ / nH} &
{$M_\mathrm{simu}$ / nH} &
{$\varepsilon_{r,\mathrm{th}}$ / \%} &
{$\varepsilon_{r,\mathrm{simu}}$ / \%} \\

\midrule

120.10 & 1.49 & 118.20 & 118.22 & 1.58 & 1.57 \\
 95.71 & 1.30 &  92.74 &  92.78 & 3.10 & 3.06 \\
 68.85 & 0.86 &  67.09 &  67.12 & 2.56 & 2.51 \\
 56.34 & 6.31 &  50.03 &  50.04 & 11.20 & 11.18 \\
 33.37 & 3.63 &  29.53 &  29.53 & 11.51 & 11.51 \\

\bottomrule

\end{tabular}

\smallskip
\raggedright
\footnotesize
$\varepsilon_{r} =
\lvert M_\mathrm{val} - M_\mathrm{exp} \rvert / M_\mathrm{exp} \times 100\,\%$.

\end{table}

\pagebreak

% ============================================================
% SECTION : PROPAGATION DES INCERTITUDES
% ============================================================

\section{Propagation des incertitudes}

Les incertitudes sont propagées au premier ordre selon la formule générale
pour une fonction $f(x_1, x_2, \ldots)$ :
\begin{equation}
  \delta f = \sqrt{\sum_i \left(\frac{\partial f}{\partial x_i}\,\delta x_i\right)^2}
  \label{eq:prop_generale}
\end{equation}

%------------------------------------------------------------
\subsection{1 bobine --- position \texorpdfstring{$z$}{z}}

La position axiale corrigée est :
\begin{equation}
  z = z_\text{sans offset} - z_\text{offset}
\end{equation}

Les deux grandeurs sont mesurées indépendamment avec une règle,
$\delta z_\text{sans offset} = \delta z_\text{offset} = \pm 1\,\text{mm}$.
La propagation donne :
\begin{equation}
  \delta z
  = \sqrt{\left(\frac{\partial z}{\partial z_\text{sans offset}}\,\delta z_\text{sans offset}\right)^2
          + \left(\frac{\partial z}{\partial z_\text{offset}}\,\delta z_\text{offset}\right)^2}
  = \sqrt{(\delta z_\text{sans offset})^2 + (\delta z_\text{offset})^2}
  = \sqrt{1^2 + 1^2}
\end{equation}
\begin{equation}
  \boxed{\delta z = \sqrt{2} \approx 1{,}41\,\text{mm}}
\end{equation}

%------------------------------------------------------------
\subsection{1 bobine --- champ magnétique \texorpdfstring{$B_\text{eff}$}{Beff}}

Le champ est déduit des tensions mesurées par la sonde~:
\begin{equation}
  B_\text{eff}
  = \frac{\sqrt{U_{x,\text{eff}}^2 + U_{z,\text{eff}}^2}}
         {\pi\!\left(\dfrac{d}{2}\right)^{\!2} \cdot N_s \cdot 2\pi f}
  \label{eq:B}
\end{equation}

où $d = 13{,}5\,\text{mm}$ est le diamètre de la bobine de mesure,
$N_s = 80$ le nombre de spires de la sonde,
et $f = 900\,\text{Hz}$ la fréquence d'excitation.\\
On pose $U = \sqrt{U_{x,\text{eff}}^2 + U_{z,\text{eff}}^2}$ et
$\mathcal{D} = \pi\!\left(\tfrac{d}{2}\right)^2 N_s\,2\pi f$,
donc $B = U / \mathcal{D}$.

\medskip
\textbf{Incertitude sur $U$ :}
$\delta U_{x,\text{eff}} = \delta U_{z,\text{eff}} = \delta U_\text{mes} = \pm 0{,}01\,\text{mV}$.
\begin{equation}
  \delta U
  = \sqrt{\left(\frac{U_{x,\text{eff}}}{U}\,\delta U_\text{mes}\right)^2
          + \left(\frac{U_{z,\text{eff}}}{U}\,\delta U_\text{mes}\right)^2}
  = \frac{\delta U_\text{mes}}{U}
    \sqrt{U_{x,\text{eff}}^2 + U_{z,\text{eff}}^2}
  = \delta U_\text{mes}
  \label{eq:dU}
\end{equation}

\textbf{Incertitude sur $\mathcal{D}$ :}
Le diamètre de la sonde est mesuré avec $\delta d = \pm 0{,}5\,\text{mm}$.
\begin{equation}
  \mathcal{D} = \frac{\pi N_s \cdot 2\pi f}{4}\,d^2
  \qquad\Rightarrow\qquad
  \frac{\delta \mathcal{D}}{\mathcal{D}} = 2\,\frac{\delta d}{d}
\end{equation}

\textbf{Propagation sur $B$ :}
\begin{equation}
  \frac{\delta B}{B}
  = \sqrt{\left(\frac{\delta U}{U}\right)^2 + \left(\frac{\delta \mathcal{D}}{\mathcal{D}}\right)^2}
  = \sqrt{\left(\frac{\delta U_\text{mes}}{U}\right)^2 + \left(2\,\frac{\delta d}{d}\right)^2}
  \label{eq:dB_rel}
\end{equation}

\begin{equation}
  \boxed{
  \delta B
  = B\,\sqrt{\left(\frac{\delta U_\text{mes}}{\sqrt{U_{x,\text{eff}}^2+U_{z,\text{eff}}^2}}\right)^{\!2}
             + \left(\frac{2\,\delta d}{d}\right)^{\!2}}
  }
\end{equation}

\pagebreak

%------------------------------------------------------------
\subsection{2 bobines --- pente de \texorpdfstring{$U_i(\omega)$}{Ui(omega)}}

L'incertitude sur la pente issue de \texttt{LINEST} est un écart-type $s_b$.
Pour obtenir une incertitude élargie à 95\,\%, d'après la littérature,
on applique un facteur à $2{,}3$.\\
L'incertitude élargie sur la pente est donc $\delta(\text{pente}) = 2{,}3\,s_b$.

%------------------------------------------------------------
\subsection{2 bobines --- inductance mutuelle \texorpdfstring{$M$}{M}}

\begin{equation}
  M = \frac{\text{pente}}{N^2\,I_\text{eff}}
\end{equation}

Les sources d'incertitude sont la pente et le courant
$\delta I_\text{eff} = \pm 0{,}01\,\text{A}$ ($N = 98$ exact, incertitude nulle).

\begin{equation}
  \frac{\delta M}{M}
  = \sqrt{\left(\frac{\delta(\text{pente})}{\text{pente}}\right)^2
          + \left(\frac{\delta I_\text{eff}}{I_\text{eff}}\right)^2}
\end{equation}

\begin{equation}
  \boxed{
  \delta M
  = \frac{1}{N^2\,I_\text{eff}}
    \sqrt{\left(\delta(\text{pente})\right)^2
          + \left(\frac{\text{pente}}{I_\text{eff}}\,\delta I_\text{eff}\right)^2}
  }
\end{equation}

\textit{Évaluation numérique pour $d = 30\,\text{mm}$}
(pente $= 1{,}038 \times 10^{-3}\,\text{V}{\cdot}\text{s/rad}$,
$\delta(\text{pente}) \approx \delta a_\text{stat}$ à estimer depuis les résidus,
$I_\text{eff} = 0{,}9\,\text{A}$)~:

\begin{equation}
  \frac{\delta I_\text{eff}}{I_\text{eff}} = \frac{0{,}01}{0{,}9} = 1{,}1\,\%
\end{equation}

Le terme dominant est en général $\delta a_\text{stat}$ ;
la contribution de $\delta I_\text{eff}$ est inférieure à $\sim 1\,\%$.

\end{document}